\documentclass{article}
\usepackage{amsfonts}
\usepackage{amsmath}
\usepackage{amsthm}
\usepackage{amssymb}
\usepackage{mathtools}% http://ctan.org/pkg/mathtools
\usepackage{hyperref}
\usepackage{caption}
\usepackage{nameref}
\usepackage[title,titletoc]{appendix}
\usepackage[margin=0.75in]{geometry}
\usepackage[fontsize=14pt]{scrextend}
\theoremstyle{definition}
\newtheorem{definition}{Definition}[subsection]

\title{Unlimited Theorems}
\date{}

\begin{document}
\maketitle
\begin{abstract}
    \noindent
    Theorems and results have been organized into sections and subsections. Definitions, if needed, are provided at the beginning of each section and subsection. A subsection has final index 0 if and only if it is a subsection of definitions. For example, subsections 5.0, 3.3.0, and 0.0 are definition subsections, but 2.3 and 1.4.1 are not.
    \newline
    \newline
    \noindent    
    If a term is defined more than once in this document, the definition provided at most specific layer will apply. For example, if section 5.0 defines ``ring" to not necessarily be commutative, but 5.3.0 declares a ``ring" to be commutative with respect to multiplication, then in all theorems 5.3.X, ``ring" will mean commutative ring, unless further overridden. Mentions of rings in 5.2, on the other hand, will not imply commutativity unless the definition is overridden.
    \newline
    \newline
    \noindent
    If a term is used but not defined within a section or subsection, the definition that is ``outer-most'' applies. For example, if the term ``homomorphism'' is encountered in section 4.1, but is defined at 2.1.0, 3.0, and 5.2.3.0, then the definition in 3.0 applies. If there is no single ``outer-most'' definition, then the earliest of the tied definitions applies. For example, if the same term is defined in 2.0, 3.0, and 4.1.0, then its use in section 5 will correspond to the definition in 2.0.
    \newline
    \newline
    \noindent
    An appendix is provided at the end of the document listing every term, along with references to each of its definitions, in alphabetical order.

\end{abstract}
\newpage

\tableofcontents





\newpage

\section{Logic}

\setcounter{subsection}{-1}
\subsection{Definitions}
\begin{definition}[Set]
    A \textbf{set} is an unordered collection of items; the items are usually called elements.
\end{definition}
\begin{definition}[Subset]
    Let $X$ be a set. A set $A$ which satisfies $x\in A \implies x\in X$ is called a \textbf{subset} of $X$. We write $A\subset X$ to mean $A$ is a subset of $X$.
\end{definition}
\begin{definition}[Power Set] Let $X$ be a set. The \textbf{power set} of $X$, denoted $\mathcal{P}(X)$ is defined to be the set of all subsets of $X$.
\end{definition}

\newpage


\section{Topology}
\setcounter{subsection}{-1}
\subsection{Definitions}

    \begin{definition}[Topology]\label{sub:topology} Let $X$ be a set. A subset $\tau$ of the power set $\mathcal{P}(X)$ is called a \textbf{topology} on $X$ if the following are true:
        \begin{itemize}
            \item $\phi,X \in \tau$
            \item Any arbitrary union of elements of $\tau$ belongs to $\tau$
            \item Any finite intersection of elements of $\tau$ belongs to $\tau$
        \end{itemize}
    \end{definition}

\subsection{Point Set Topology}
\ref{sub:topology}


\newpage
\renewcommand{\appendixpagename}{Appendix}
\appendixpage
\appendix
\section{Definition Index}
\section{Notation}

\end{document}
